\section{Kth Element of Two Sorted Arrays}
\label{sec:bs-kth-element}

\problemstatement{We are given two sorted arrays $\textit{nums1}$ (size
$n_1$) and $\textit{nums2}$ (size $n_2$), and an integer $k$. We must find
the $k$-th smallest element (1-indexed) of the combined sorted array,
without merging, in $O(\log(\min(n_1,n_2)))$ time.}

\begin{patternbox}
This is the direct generalization of
Section~\ref{sec:bs-median-two-arrays}: the median is simply the special
case $k = \lceil (n_1+n_2)/2 \rceil$ (for odd total length) or the average
of the $k$-th and $(k{+}1)$-th elements (for even total length). The same
``binary search on a partition split between two arrays'' technique
applies, with the partition sized to put exactly $k$ elements on the left
instead of exactly half.
\end{patternbox}

\subsection*{Understanding the problem carefully}

We again look for a split $(\textit{cut1}, \textit{cut2})$ with
$\textit{cut1} + \textit{cut2} = k$, $0 \le \textit{cut1} \le n_1$, and
$0 \le \textit{cut2} \le n_2$, such that the left side (the first
$\textit{cut1}$ elements of $\textit{nums1}$ and first $\textit{cut2}$
elements of $\textit{nums2}$) contains exactly the $k$ smallest elements
overall. The same four boundary values $l_1, r_1, l_2, r_2$ and the same
validity condition $l_1 \le r_2$ and $l_2 \le r_1$ from
Section~\ref{sec:bs-median-two-arrays} apply unchanged.

\begin{ideabox}[Candidate Space and Predicate]
The candidate space is $\textit{cut1} \in [\max(0, k-n_2), \min(k, n_1)]$
(the lower bound ensures $\textit{cut2} = k - \textit{cut1} \le n_2$; the
upper bound ensures $\textit{cut1} \le n_1$ and $\textit{cut2} \ge 0$).
The predicate is the same validity check as
Section~\ref{sec:bs-median-two-arrays}: $P(\textit{cut1}) = (l_1 \le r_2
\text{ and } l_2 \le r_1)$.
\end{ideabox}

\subsection*{Why only the target split size changes}

Nothing about the underlying correctness argument changes from
Section~\ref{sec:bs-median-two-arrays} -- we are still looking for the
unique partition where the left side's largest element does not exceed
the right side's smallest element, guaranteeing the left side is exactly
the $k$ smallest elements of the combined array. Once a valid partition is
found, the $k$-th smallest element is simply $\max(l_1, l_2)$ -- the
largest element among the $k$ elements placed on the left.

\subsection*{Algorithm}

\begin{enumerate}
  \item $\textit{lo} \gets \max(0, k-n_2)$,
        $\textit{hi} \gets \min(k, n_1)$.
  \item While $\textit{lo} \le \textit{hi}$:
  \begin{enumerate}
    \item $\textit{cut1} \gets \textit{lo}+(\textit{hi}-\textit{lo})/2$,
          $\textit{cut2} \gets k-\textit{cut1}$.
    \item Compute $l_1, r_1, l_2, r_2$ exactly as in
          Section~\ref{sec:bs-median-two-arrays}.
    \item If $l_1 \le r_2$ and $l_2 \le r_1$: return $\max(l_1,l_2)$.
    \item Else if $l_1 > r_2$: $\textit{hi} \gets \textit{cut1}-1$.
    \item Else: $\textit{lo} \gets \textit{cut1}+1$.
  \end{enumerate}
\end{enumerate}

\begin{algorithm}[H]
\caption{Kth Element of Two Sorted Arrays}
\begin{algorithmic}[1]
\Procedure{KthElement}{$\textit{nums1}, \textit{nums2}, k$}
  \State $\textit{lo} \gets \max(0, k-n_2)$, $\textit{hi} \gets \min(k, n_1)$
  \While{$\textit{lo} \le \textit{hi}$}
    \State $\textit{cut1} \gets \textit{lo}+(\textit{hi}-\textit{lo})/2$
    \State $\textit{cut2} \gets k - \textit{cut1}$
    \State $l_1 \gets (\textit{cut1}=0)\ ? -\infty : \textit{nums1}[\textit{cut1}-1]$
    \State $r_1 \gets (\textit{cut1}=n_1)\ ? +\infty : \textit{nums1}[\textit{cut1}]$
    \State $l_2 \gets (\textit{cut2}=0)\ ? -\infty : \textit{nums2}[\textit{cut2}-1]$
    \State $r_2 \gets (\textit{cut2}=n_2)\ ? +\infty : \textit{nums2}[\textit{cut2}]$
    \If{$l_1 \le r_2$ \textbf{and} $l_2 \le r_1$}
      \State \Return $\max(l_1, l_2)$
    \ElsIf{$l_1 > r_2$}
      \State $\textit{hi} \gets \textit{cut1} - 1$
    \Else
      \State $\textit{lo} \gets \textit{cut1} + 1$
    \EndIf
  \EndWhile
\EndProcedure
\end{algorithmic}
\end{algorithm}

\subsection*{Dry run}

\dryrun{Take $\textit{nums1} = [2,3,6,7,9]$, $\textit{nums2} = [1,4,8,10]$,
$k = 5$.}

$n_1=5$, $n_2=4$. $\textit{lo}=\max(0,5-4)=1$,
$\textit{hi}=\min(5,5)=5$.

\begin{itemize}
  \item $\textit{lo}=1,\textit{hi}=5$: $\textit{cut1}=3$,
        $\textit{cut2}=5-3=2$. $l_1=\textit{nums1}[2]=6$,
        $r_1=\textit{nums1}[3]=7$. $l_2=\textit{nums2}[1]=4$,
        $r_2=\textit{nums2}[2]=8$. Check: $l_1=6 \le r_2=8$? Yes.
        $l_2=4 \le r_1=7$? Yes. Valid partition.
\end{itemize}

The answer is $\max(l_1,l_2) = \max(6,4) = 6$. Checking by hand: the
merged array is $[1,2,3,4,6,7,8,9,10]$, whose $5$-th smallest element
(1-indexed) is indeed $6$.

\subsection*{C++ implementation}

\begin{lstlisting}
#include <bits/stdc++.h>
using namespace std;

int kthElement(vector<int>& nums1, vector<int>& nums2, int k) {
    int n1 = nums1.size(), n2 = nums2.size();

    int lo = max(0, k - n2);
    int hi = min(k, n1);

    while (lo <= hi) {
        int cut1 = lo + (hi - lo) / 2;
        int cut2 = k - cut1;

        int l1 = (cut1 == 0) ? INT_MIN : nums1[cut1 - 1];
        int r1 = (cut1 == n1) ? INT_MAX : nums1[cut1];
        int l2 = (cut2 == 0) ? INT_MIN : nums2[cut2 - 1];
        int r2 = (cut2 == n2) ? INT_MAX : nums2[cut2];

        if (l1 <= r2 && l2 <= r1) {
            return max(l1, l2);
        } else if (l1 > r2) {
            hi = cut1 - 1;
        } else {
            lo = cut1 + 1;
        }
    }
    return -1; // unreachable for valid input with 1 <= k <= n1+n2
}

int main() {
    vector<int> nums1 = {2, 3, 6, 7, 9};
    vector<int> nums2 = {1, 4, 8, 10};
    cout << "5th smallest element: "
         << kthElement(nums1, nums2, 5) << endl;
    return 0;
}
\end{lstlisting}

\begin{complexitybox}
\textbf{Time Complexity.} The binary search range has width at most
$\min(k, n_1) - \max(0, k-n_2) \le \min(n_1, n_2)$, giving
\[
  O(\log(\min(n_1, n_2))).
\]
\textbf{Space Complexity.} $O(1)$ auxiliary space.
\end{complexitybox}

\begin{takeawaybox}
Median-of-two-sorted-arrays and $k$-th-element-of-two-sorted-arrays are the
same partition-search technique with only the target split size changed
from ``half'' to ``$k$.'' Whenever you solve one binary-search-on-a-
partition problem, immediately check whether nearby sheet problems are
just different target split sizes of the same idea.
\end{takeawaybox}
